\section{Round-twenty-three reconstruction: defect coordinates and the correct collision-action Hessian}
\label{sec:r23-b3}

The former second-variation formula assigned positive cost to motion along the
zero-cost manifold $\Gamma=A_f$.  The active paper uses the collision defect
\[
 h=\dot\Gamma-D A_f[u],                                  \tag{B3.1}
\]
not the separate logarithmic changes of $A_f$ and $\Gamma$, as the normal
coordinate.  At an unforced Boltzmann path the collision part of the Hessian is
exactly $\frac12\int h^2/A_f$ and therefore vanishes when
$\dot\Gamma=D A_f[u]$.

\subsection{Regular reference path and global collision geometry}

Let $f$ be a positive classical Boltzmann solution on $[0,T]$ satisfying, for
some Maxwellians $M_-,M_+$ and constants $0<c<C<\infty$,
\[
 cM_-(v)\le f(t,x,v)\le CM_+(v),
 \qquad
 \sum_{|\alpha|\le4}|\partial_{x,v}^\alpha f|
 \le C M_+(v)^{1/2}.                                    \tag{B3.2}
\]
This global comparison, obtained on the regular preparation class propagated
in B2, replaces the insufficient compact-by-compact positivity of the previous
version.  Put
\[
 A_f=f f_*[(v-v_*)\cdot\omega]_+
 \,dt\,dx\,dv\,dv_*\,d\omega.                           \tag{B3.3}
\]
For a density tangent $u$, differentiation gives
\[
 D A_f[u]=A_f\left({u\over f}+{u_*\over f_*}\right).     \tag{B3.4}
\]
Let $\mathcal C$ send a signed collision measure to its gain-minus-loss
velocity marginal.  The linearized balance relation in defect coordinates is
\[
 \partial_tu+v\cdot\nabla_xu
 =D Q_f[u]+\mathcal C h,                                  \tag{B3.5}
\]
where $DQ_f[u]=\mathcal C(D A_f[u])$.

For an adjoint test $p$, write
$\Delta p=p'+p_*'-p-p_*$.  The microscopic Dirichlet form is
\[
 \mathcal D_f(p)=\frac14\int A_f(\Delta p)^2.             \tag{B3.6}
\]
By (B3.2), it is globally comparable with the Maxwellian cutoff collision
form.  Its kernel is precisely the five collision invariants
$1,v_1,v_2,v_3,|v|^2$.

\begin{lemma}[Global microscopic gap]
\label{lem:r23-b3-gap}
Let $\Pi$ be the orthogonal projection onto the collision invariants.  In the
weighted space $L^2(M_+^{-1/2})$,
\[
 \|(I-\Pi)p\|_{L^2(M_+)}^2
 \le C\mathcal D_f(p).                                   \tag{B3.7}
\]
The constant is uniform on the regular reference class (B3.2).
\end{lemma}

\begin{proof}
The cutoff Maxwellian collision form has a spectral gap on the orthogonal
complement of its five-dimensional invariant space.  Global two-sided
comparison (B3.2) transfers both the measure and the form norms with fixed
constants.  Truncation in velocity and monotone convergence justify the
comparison on the full weighted space.  No velocity $H^1$ gain is asserted.
\end{proof}

\subsection{Transport--collision observability and closed range}

Define the balance operator
\[
 \mathfrak B(u,h)=
 \bigl(\partial_tu+v\cdot\nabla_xu-DQ_f[u]-\mathcal C h,
       u(0)\bigr)                                        \tag{B3.8}
\]
on the graph space obtained by completing smooth specular functions under
\[
 \|u(0)\|_{\mathcal H_0}^2+
 \int_0^T\|(I-\Pi)u_t\|_{L^2(M_+^{-1})}^2dt+
 \int {h^2\over A_f}.                                   \tag{B3.9}
\]
The macroscopic coefficients $\Pi u$ are controlled by the transport equation
and conservation constraints.  Fourier transform in $x$, apply the
microscopic estimate (B3.7), and for each nonzero spatial mode use the
Kawashima cross term between $ik\cdot v\Pi u$ and $(I-\Pi)u$.  The zero mode is
fixed by mass, momentum, and energy.  Integration in time gives
\[
 \|u\|_{L^2_t\mathcal H_x}^2
 \le C\left(\|u(0)\|_{\mathcal H_0}^2+
            \int {h^2\over A_f}+
            \|\mathfrak B(u,h)_1\|_{\mathcal Y}^2\right).              \tag{B3.10}
\]
Coefficient variation in time is controlled by (B3.2) and absorbed after a
finite time partition.

\begin{theorem}[Closed balance range modulo invariants]
\label{thm:r23-b3-range}
After imposing the five conservation constraints, $\mathfrak B$ has closed
range and admits a bounded right inverse on that range.  Its null directions
are exactly the solutions of the homogeneous linearized Boltzmann equation
with allowed initial tangent.  No assertion is made that an unrelated
cokernel is finite dimensional; the only finite-dimensional space used is the
explicit collision-invariant space.
\end{theorem}

\begin{proof}
Estimate (B3.10) is the graph coercivity estimate modulo the stated null
space.  A Cauchy sequence of images therefore has representatives Cauchy
modulo that null space; fixing the orthogonal gauge gives convergence in the
graph norm and proves closed range.  The inverse on the orthogonal complement
is bounded by the closed graph theorem.  Testing the homogeneous equation
against the adjoint and using (B3.7) identifies the null directions.
\end{proof}

\subsection{A nuclear process space and stopping-time cumulants}

Choose Fourier modes in $x$, Hermite modes in $v$, and spherical harmonics in
$\omega$.  Let $\mathcal S_r$ be the Hilbert space with coefficient weight
$(1+|k|+|n|+|\ell|)^{r}$ and a Maxwellian velocity factor.  For
$r_{j+1}-r_j$ larger than half the total mode dimension, the embeddings
\[
 \mathcal S_{r_2}\hookrightarrow\mathcal S_{r_1}
 \hookrightarrow\mathcal S_{r_0}                         \tag{B3.11}
\]
are Hilbert--Schmidt.  The projective limit is nuclear and its dual carries
the joint density/contact fluctuation process.

B2 gives localized connected cumulants for deterministic intervals.  To pass
to stopping times, let $\tau_\varepsilon$ be a stopping time for the full
microscopic filtration and condition on the complete phase point
$\Phi_{\tau_\varepsilon}^{\varepsilon,N}\mathbf X$.  Deterministic
hard-sphere flow has the exact restart property on this full state.  B2's
history norm is uniform over the prepared regular set, so its connected
expansion applied after the restart gives, for $\theta\le\delta$,
\[
 E\left[
 \|Z^\varepsilon_{\tau_\varepsilon+\theta}
       -Z^\varepsilon_{\tau_\varepsilon}\|_{\mathcal S_{r_2}'}^2
 \mid\mathcal F_{\tau_\varepsilon}\right]
 \le C\delta+o_\varepsilon(1).                           \tag{B3.12}
\]
The maximum jump is $o(1)$ by the collision-history factorial bound and the
grazing truncation.  Thus Aldous tightness follows from a conditional restart
estimate, not from independent increments or an unconditional deterministic
interval bound.

\begin{theorem}[Joint Gaussian process]
\label{thm:r23-b3-clt}
The centered density/contact fluctuations converge in
$D([0,T];\mathcal S_{r_2}')$ to the continuous Gaussian solution of the
linearized balance equation driven by collision noise with covariance
\[
 E[dW_c(\psi)dW_c(\phi)]
 =\int A_f\,\Delta\psi\,\Delta\phi.                      \tag{B3.13}
\]
Microcanonical constraints are imposed by the B1 Schur projection of the
initial covariance.
\end{theorem}

\begin{proof}
B2's second connected cumulant converges to (B3.13), while all connected
cumulants of order at least three vanish after central-limit scaling by the
factorial hierarchy estimate.  Finite-dimensional convergence follows by the
cumulant expansion.  Estimate (B3.12), the vanishing jumps, and the
Hilbert--Schmidt tail of (B3.11) give process tightness.  The martingale
problem determined by (B3.5) and (B3.13) has a unique Gaussian solution.
Exact-number conditioning acts on the joint initial Gaussian and therefore
uses the Schur complement proved in B1.
\end{proof}

\subsection{Exact second epi-derivative}

The collision entropy is
\[
 \mathcal J(f,\Gamma)=\int A_f\ell(q),
 \qquad q={d\Gamma\over dA_f},
 \qquad \ell(q)=q\log q-q+1.                             \tag{B3.14}
\]
For comparison with the referee's calculation, suppose
$A_{f_\eta}=e^{\eta a+O(\eta^2)}A_f$ and
$q_\eta=qe^{\eta k+O(\eta^2)}$.  Direct differentiation gives
\[
 {d^2\over d\eta^2}\bigg|_0 e^{\eta a}\ell(qe^{\eta k})
 =\ell(q)a^2+2q\log(q)ak+q(1+\log q)k^2                 \tag{B3.15}
\]
before the declared second-order chart terms.  The former expression
$q|k+a|^2$ is not used.

At the zero-cost reference $q=1$, use the normal coordinate (B3.1).  For a
smooth tangent $(u,h)$ satisfying (B3.5), choose the positive chart
\[
 f_\eta=f\exp\{\eta u/f-\psi_\eta\},
 \qquad
 \Gamma_\eta=A_{f_\eta}
 \exp\left\{\eta {h\over A_f}+r_\eta\right\}.            \tag{B3.16}
\]
The scalar/vector normalizers $\psi_\eta$ preserve mass, momentum, and energy.
The balance residual of the first-order chart is $O(\eta^2)$.  Apply the
bounded right inverse from Theorem~\ref{thm:r23-b3-range} to choose
$r_\eta=O(\eta^2)$ in the graph norm so that the nonlinear balance equation
holds exactly.  Truncation followed by diagonal removal treats unbounded
$h/A_f$ while preserving the limiting cost.

\begin{theorem}[Correct Mosco second epi-derivative]
\label{thm:r23-b3-mosco}
Let $I$ be the B2 action including the prepared initial entropy.  At a
zero-cost Boltzmann path $(f,A_f)$,
\[
 d_e^2I(f,A_f)[u,h]
 =\frac12\|u_0\|_{\Sigma_0^{-1}}^2
  +\frac12\int {h^2\over A_f},                           \tag{B3.17}
\]
subject to (B3.5) and the microcanonical tangent constraints, and it is
$+\infty$ otherwise.  The forms Mosco-converge on the nuclear tangent triple,
and the inverse on the closed range is the covariance of
Theorem~\ref{thm:r23-b3-clt}.
\end{theorem}

\begin{proof}
For the lower bound, write
$d\Gamma_\eta/dA_{f_\eta}=1+\eta h/A_f+o(\eta)$ weakly on finite-action
sublevels.  Since
$\ell(1+x)=x^2/2+o(x^2)$ and the entropy integrands are uniformly integrable,
convex lower semicontinuity gives
$\liminf I/\eta^2\ge\frac12\int h^2/A_f$, together with the prepared initial
quadratic form.  Passing the balance equation to first order gives (B3.5).
For the upper bound, use the exactly balanced positive chart (B3.16); Taylor's
formula and the $O(\eta^2)$ right-inverse correction give equality in the
limit.  Graph coercivity yields equicoercivity and hence Mosco convergence.
The Gaussian Cameron--Martin norm associated with (B3.13) is precisely the
minimum of (B3.17) over collision noises producing a given density tangent,
which identifies the inverse form with the covariance.
\end{proof}

If $h=0$, then $\Gamma_\eta=A_{f_\eta}$ to first order and the collision part
of (B3.17) is zero.  This includes the referee's test $q=1$, $k=0$, and
arbitrary variation of $A_f$ along the zero-cost manifold.  Thus the explicit
counterexample is now a passed regression identity rather than an exception.

\subsection{Literature boundary}

The fluctuation covariance agrees with the collision noise appearing in the
hard-sphere fluctuation theory of
Bodineau--Gallagher--Saint-Raymond--Simonella \cite{BGSS2023}.  The present
paper's additional tasks are the stopped-history tightness, the closed
balance range, and the defect-coordinate second epi-derivative.  None is
inferred solely from the existence of deterministic-time cumulants.
